% !PRINTED 594
\weberpage{594}
%<<CONT
\gesperrt{Fermat's} theorem for rational numbers as well, one recognises that the congruences (2) and (5) are \gesperrt{equivalent in meaning} to the formulae valid for \gesperrt{every number} $\omega$ in $\mathfrak{o}$

\weq{6}{\omega\mkern1mu|\mkern1mu\chi \equiv \omega,\quad \omega^P \equiv \omega\mkern1mu|\mkern1mu\psi_0\ (\mathrm{mod}\ \mathfrak{p}).}
From the second congruence (6) it follows that, if $\omega$ is divisible by $\mathfrak{p}$, then $\omega\mkern1mu|\mkern1mu\psi_0$ must always be divisible by $\mathfrak{p}$ as well. Consequently $\mathfrak{p}$ is divisible by $\mathfrak{p}\mkern1mu|\mkern1mu\psi_0$, and, being a prime ideal, $= \mathfrak{p}\mkern1mu|\mkern1mu\psi_0$. From this it follows that $\psi_0$ is contained in the group $\Psi$.

We now understand by $\gamma$ a primitive root of $\mathfrak{p}$, and apply the formula proved at the end of \S.\,161:
% [sic: The source repeats $\gamma^{P^f-1}$ in the last factor, as on printed page 592; the P-power orbit instead requires $\gamma^{P^{f-1}}$. The printed expression is retained.]
\weq{7}{(t-\gamma)(t-\gamma^P)\ldots(t-\gamma^{P^f-1}) \equiv F(t)\ (\mathrm{mod}\ \mathfrak{p}),}
in which $F(t)$ is an integral function of $t$ in $R$.

From this it follows that
\[ F(\gamma) \equiv 0\ (\mathrm{mod}\ \mathfrak{p}), \]
and if now $\psi$ is any substitution from $\Psi$:
\[ F(\gamma\mkern1mu|\mkern1mu\psi) \equiv 0\ (\mathrm{mod}\ \mathfrak{p}). \]

But from this it results that $\gamma\mkern1mu|\mkern1mu\psi$ must be congruent, with respect to $\mathfrak{p}$, to one of the powers of $\gamma$ occurring in (7), and so, say,
\weq{8}{\gamma^{P^\nu} \equiv \gamma\mkern1mu|\mkern1mu\psi\ (\mathrm{mod}\ \mathfrak{p}).}
Now every number $\omega$ in $\mathfrak{o}$ not divisible by $\mathfrak{p}$ is congruent to a power of $\gamma$, and if one therefore raises (8) to this power, it follows that
\weq{9}{\omega^{P^\nu} \equiv \omega\mkern1mu|\mkern1mu\psi\ (\mathrm{mod}\ \mathfrak{p}),}
and this congruence evidently holds even when $\omega$ is divisible by $\mathfrak{p}$, and hence for \gesperrt{all numbers} in $\mathfrak{o}$.

If we now apply the second of the congruences (6) $\nu$ times in succession, it follows that
\weq{10}{\omega^{P^\nu} \equiv \omega\mkern1mu|\mkern1mu\psi_0^\nu\ (\mathrm{mod}\ \mathfrak{p}),}
hence
\weq{11}{\omega\mkern1mu|\mkern1mu\psi \equiv \omega\mkern1mu|\mkern1mu\psi_0^\nu\ (\mathrm{mod}\ \mathfrak{p}),}
and if we replace $\omega$ by $\omega\mkern1mu|\mkern1mu\psi^{-1}$, or apply the substitution $\psi^{-1}$ to (11):
\weq{12}{\omega \equiv \omega\mkern1mu|\mkern1mu\psi^{-1}\psi_0^\nu \equiv \omega\mkern1mu|\mkern1mu\psi_0^\nu\psi^{-1}\ (\mathrm{mod}\ \mathfrak{p}).}
Thus both $\psi^{-1}\psi_0^\nu$ and $\psi_0^\nu\psi^{-1}$ are contained in $X$,
%CONT>>
